\chapter{酒精、香气与酸味}

\begin{kaichang}
请想象你的面前摆着三样东西：一瓶白醋、一根熟透了的香蕉、一小瓶医用消毒酒精。

先别急着说它们“毫无关系”。来做几个思想实验：一瓶葡萄酒开瓶后放上几个星期，常常会变酸，隐约有了醋味；一串香蕉熟过头，会散发出一丝若有若无的酒味；而用粮食、水果发酵，既能造出酒，也能造出醋。这三样东西之间，仿佛藏着一条隐形的流水线，东西放上去，就会一个接一个地变成另一样。

猜一猜：这条流水线是什么？酒、香气、酸味这三种完全不同的体验，怎么会是同一条线上的三个车站？这一章结束时，你会用一张只有四个“零件”的图纸，把整条流水线画出来。
\end{kaichang}

\section{三瓶液体，一条隐形流水线}

先把镜头停在眼睛、鼻子和皮肤能直接报告的层面上。

白醋倒在杯子里，是无色透明的液体，凑近了有一股尖锐刺鼻的酸味，蘸一点在舌尖上，酸得让人眯眼。它能和水以任何比例混在一起，永远不分层。

医用酒精也是无色透明的，但脾气完全不同：倒在手背上，几秒钟就挥发得无影无踪，留下一片凉意——挥发要带走热量（第 4 章讲过蒸发吸热）。它也能和水任意混合，而且很容易点燃，烧起来是安静的淡蓝色火焰。

香蕉的香气看不见摸不着，但它会自己钻进你的鼻子——这说明有某种分子源源不断地从果肉里跑出来，飘过空气，落进你鼻腔里的感受器上。能闻到的，必定是容易挥发的分子。

三种物质，三种性格。但它们有一个不声不响的共同点：凑近看分子式，都由碳、氢、氧三种原子组成，而且氧原子都不多，常常只有一两个。氧原子这么少，为什么存在感这么强？

\section{氧原子是碳骨架上的“接口”}

第 36 章说过，碳原子是搭骨架的好手：四个键，能连成链、连成环。可是纯碳氢骨架——比如第 39 章的烷烃——性格相当沉闷：不溶于水，不爱反应，烧掉是它最激动的时刻。

氧原子一插进来，局面立刻改观。氧是周期表里出了名的“电子磁铁”（第 11 章的电负性），它一上车，就把附近的电子往自己这边拉，让原本均匀的骨架出现“局部带电”的热点。电子分布不均匀的地方，就是水分子喜欢停靠的地方，也是别的分子发起攻击的地方——第 37 章把这样的位置叫做\textbf{官能团}，分子的“反应接口”。

这一章的主角只有四种接口，先认个脸熟：

\begin{itemize}[itemsep=2pt]
\item \textbf{羟基}：一个氧连一个氢（$-\mathrm{OH}$），挂在碳链上，这个分子就叫\textbf{醇}。乙醇、甲醇都是醇。
\item \textbf{羰基}：一个碳和一个氧用双键相连（$\mathrm{C{=}O}$）。羰基在碳链的头上，叫\textbf{醛}；羰基在碳链中间，叫\textbf{酮}。乙醛是醛，洗甲水里的丙酮是酮。
\item \textbf{羧基}：羰基和羟基长在同一个碳上（$-\mathrm{COOH}$），这个组合有了新名字和新性格，叫\textbf{羧酸}。醋里的乙酸就是羧酸。
\item \textbf{酯基}：羧基里的氢换成一段碳链（$-\mathrm{COO}-$），就成了\textbf{酯}。香蕉、菠萝的香气分子，多半是酯。
\end{itemize}

看明白了吗？这四种接口是同一个主题的四次变奏：羰基、羟基两种“零件”拼来拼去。开场那条隐形流水线，其实就是氧原子在碳骨架上换装的过程——下面一步步看。

\section{氢键：水溶性背后的“握手”}

先解决一个现象层的谜题：为什么酒精和醋都能和水任意混合，而第 39 章的汽油倒进水里却死活分层？

第 22、23 章讲过水分子的绝活：氢键。水分子之间靠氢键彼此拉扯，织成一张黏黏的网。一个外来分子想溶进水里，就得能跟这张网“搭上话”——自己也得会形成氢键，否则就会被水网排挤出去，聚成另外一层。

羟基恰好是氢键的“标准接口”：氧上带着微微的负电，氢上带着微微的正电，和水分子握起手来严丝合缝。所以乙醇分子混进水里，就像一滴水混进另一滴水。乙酸有羧基，同样会握手，同样任意互溶。而汽油里的烷烃分子浑身找不到这样的接口，只能被水网拒之门外。

\begin{qiaoliang}
羟基让醇“会握手”，但碳链部分仍然是“不合群”的。于是水溶性成了一场拔河：羟基往水里拉，碳链往油里拉。谁赢，看身材。甲醇（\chem{CH_3OH}）、乙醇（\chem{CH_3CH_2OH}）碳链很短，羟基说了算，与水任意互溶。丁醇有四个碳，在 $20\,^{\circ}\mathrm{C}$ 时每 100 克水大约只能溶 8 克。到了辛醇（八个碳），每 100 克水只能溶约 0.05 克——羟基的声音已经被长长的碳链淹没了。同一条规律还能解释沸点：乙醇的沸点是 $78\,^{\circ}\mathrm{C}$，而把它的羟基去掉、只留碳氢骨架的乙烷，沸点低到 $-89\,^{\circ}\mathrm{C}$。相差将近 170 度！这 170 度就是氢键这张“分子网”的身价：要烧开乙醇，得额外提供能量，把分子之间一双双握着的“手”掰开。
\end{qiaoliang}

\begin{shiyan}
\textbf{哪瓶液体和水合得来？}

材料：三个透明玻璃杯、清水、少量食用油、白醋、少量白酒或医用酒精、一根筷子。

步骤：三个杯子各装半杯水。往第一杯滴入几滴食用油，第二杯倒入一勺白醋，第三杯倒入一勺酒精。用筷子分别搅动十下，然后静置五分钟，从侧面观察。

观察点：哪杯出现了清晰的分层？哪杯完全均匀？对着光看，均匀的杯子是真透明还是有淡淡的浑浊？（猜一猜：如果把酒精换成同样多的汽油会怎样？第 39 章的读者应该已经有答案了。）

安全提示：只用少量，不用嘴尝任何实验液体；酒精务必远离明火和灶台；实验结束后把液体倒进下水道并用清水冲走。
\end{shiyan}

\section{从酒精到醋：一条氧化链}

现在把镜头推到符号层，把开场的流水线画出来。酵母菌把糖变成乙醇（这个魔术第 52 章会拆给你看），而乙醇在氧气和某些细菌面前并不安稳，它会一级一级地被氧化：

\[\chem{CH_3CH_2OH} \rightarrow \chem{CH_3CHO} \rightarrow \chem{CH_3COOH}\]

乙醇，到乙醛，再到乙酸。数一数氧原子：1 个、1 个、2 个；再看接口：羟基变成了羰基（醛），羰基又“吃进”一个氧变成羧基。所谓“酒放久了变酸”，就是空气中的醋酸菌把这条链走完了——你厨房里那瓶醋，和酒杯里的酒，确实隔着两站地。

第 33 章说过，氧化的本质是失去电子，跟“碰到氧气”没有必然关系。这条链上每一步的真正含义是：碳原子身边的电子被氧抢走了一部分，碳的“氧化态”一级一级升高。身体里的酶和醋酸菌做的，是同一种电子搬运工作，只是一个在细胞里、一个在醋缸里。

顺便认识一位危险的亲戚：甲醇（\chem{CH_3OH}），比乙醇少一个碳。它自己也走这条氧化链，产物却是甲醛（\chem{HCHO}）和甲酸（\chem{HCOOH}）——甲醛会破坏细胞里的蛋白质，甲酸堆积会扰乱血液的酸碱平衡（第 32 章），首当其冲受害的是视神经。工业酒精里常混有甲醇，这就是“假酒致人失明”的化学根源：不是酒“烈”，而是氧化链的另一条岔路通向了毒物。

\section{酒精进入身体之后}

乙醇这个分子，身材小，又会和水握手、又不太嫌弃油脂（它有短短两个碳的“油尾巴”），于是它在身体里几乎畅行无阻：从胃和小肠钻进血液，穿过细胞膜，几分钟内就能抵达大脑。在大脑里，它会干扰神经细胞之间传递信号的分子开关，让某些“减速”信号变强、“加速”信号变弱——这就是反应变慢、走路打晃的分子起点。这部分属于生理学，这里不展开；我们要跟紧的是化学：身体怎样把这个外来分子拆掉。

在讲拆解之前，先解开消毒酒精的一个谜：药店卖的消毒酒精是 75\% 左右的，为什么不用 100\% 的纯酒精——难道浓度越高杀菌越弱？还真有点这个意思。酒精杀菌靠的是钻进微生物体内、让蛋白质变性散架。纯酒精太“急”，一接触细菌表面就让表层蛋白质瞬间凝固，等于给细菌现场浇了一层硬壳，里面的酒精反而渗不进去。兑了水的酒精渗透得从容，能钻进内部再动手。这个细节值得记住，因为它打破了“越浓越好、越纯越强”的直觉——化学里，配比比纯度常常是更要紧的变量。

拆酒精的主战场在肝脏，用的正是上面那条氧化链，只不过搬运电子的不是醋酸菌，而是两种酶（酶的本质第 49 章细讲）：

\begin{itemize}[itemsep=2pt]
\item 第一种酶（乙醇脱氢酶）把乙醇氧化成乙醛；
\item 第二种酶（乙醛脱氢酶）把乙醛氧化成乙酸；
\item 乙酸随后进入全身的能量代谢，被拆成二氧化碳和水，顺便放出能量（第 53 章）。
\end{itemize}

整条链上最扎眼的是中间产物乙醛。乙醛的毒性比乙醇强得多：它能让血管扩张（于是脸红、脖子红），能引起头痛、心悸、恶心，还会直接攻击 DNA。国际癌症研究机构把酒精饮料本身、以及随饮酒进入体内的乙醛，都列入了“有充分证据对人类致癌”的 1 类清单。请注意这句话的准确含义：它说的是“证据等级充分”，不是说“碰一滴就致癌”——风险随剂量上升，这是毒理学的第一条铁律（还记得第 1 章那句“剂量决定毒性”吗）。

第二条铁律是\textbf{个体差异}。第一步酶人人都够用，瓶颈在第二步：大约有三分之一的东亚人携带一种效率大打折扣的乙醛脱氢酶（下一节的证据故事就讲它）。对这些人来说，乙醛在体内停留的时间更长、浓度更高——同样一杯酒，他们承受的乙醛冲击比别人大得多。换句话说，\textbf{不存在对所有人都一样的“安全剂量”}；而且对大脑仍在发育的未成年人，酒精的干扰本身就没有“安全”可言。本章只讲化学和证据，不提供任何关于饮酒的建议。

\begin{qiaoliang}
做一次数量级估算。一罐 330 毫升、酒精度 5\% 的啤酒，含乙醇约 $330\times 5\% \approx 16.5$ 毫升。乙醇密度约 \ud{0.79}{g/mL}，折合约 13 克；除以摩尔质量 \ud{46}{g/mol}，约 0.28 摩尔——也就是大约 $1.7\times 10^{23}$ 个乙醇分子。一个成年人肝脏平均每小时大约处理 7 克乙醇（这只是平均值，个体之间能差出好几倍），照这个速度，仅这一罐酒就需要肝脏连续工作约两个小时。你可以自己算算：如果喝得更快、更多，排队的分子会在血液里积到什么程度——“醉”，本质上就是进入身体的速度远远跑赢了拆解的速度。
\end{qiaoliang}

\begin{wuqu}
“喝酒脸红的人酒量好，练练就能喝。”——恰好说反了。脸红是乙醛堆积、血管被迫扩张的报警信号，说明第二步拆解跟不上，有毒物质正在体内滞留。带着这种体质长期饮酒的人，食道癌等疾病的风险显著高于不脸红的饮酒者。所谓“练出来的酒量”，只是大脑对酒精的干扰渐渐麻木，报警声被调小了，乙醛对 DNA 的攻击却一分没停。报警器不响了，不等于火灭了。
\end{wuqu}

\section{科学家怎样破译“脸红之谜”}

\begin{zhengju}
喝酒脸红这个现象，东亚人里常见，欧洲人里罕见。表面看是个生活趣闻，科学家却从中嗅到了一个标准的研究问题：差异到底出在哪里？

假说有好几个候选：也许是饮食习惯不同，也许是“体质”这种说不清的东西，也许干脆是心理作用。20 世纪 70 年代，研究者把问题拆成可以测量的化学问题：既然乙醇是按“乙醇→乙醛→乙酸”这条链拆解的，那么脸红的人，是链上哪一环慢了？他们测定不同人肝脏中两种脱氢酶的活性，结果发现了一个清晰的图案：许多喝酒脸红的人，乙醇脱氢酶完全正常，慢的是乙醛脱氢酶——第一种酶火力全开，把乙醇迅速变成乙醛，第二种酶却跟不上，乙醛在体内堆积。这直接预言了：脸红者血液里的乙醛浓度应该更高。抽血测定，果然如此。

可酶为什么慢？20 世纪 80 年代末，答案落到了 DNA 上：乙醛脱氢酶的基因里，一个碱基的差别导致酶蛋白上第 487 号氨基酸被换掉，酶的三维结构随之变形，组装起来的“机器”松松垮垮，活性大跌。再后来，大规模的流行病学调查把基因型、饮酒习惯和食道癌发病率放在一起比对，确认了那条让人不安的关联：携带这种变异还长期饮酒的人，风险成倍上升。

这个故事漂亮在哪？每一步都有替代解释被排除：不是“体质玄学”，因为能测到具体的酶活性差异；不是“喝得急”，因为基因型能预言谁在同样饮酒量下乙醛更高；而最终的因果链——一个碱基→一个氨基酸→一台变形的酶→堆积的乙醛→升高的风险——每一环都有独立证据。它也提前剧透了第三册的核心剧情：DNA 上的微小差别，怎样通过蛋白质的形状，一路放大成肉眼可见的个体差异（第 71、72 章会彻底讲透）。
\end{zhengju}

\section{羰基：一根会“吸电子”的双键}

流水线的中间站是醛和酮，它们的共同零件是羰基：碳和氧双键相连。第 38 章讲过分析有机反应的通用心法——先看电子在哪里挤、在哪里缺。羰基就是教科书级的例子：氧把双键里的电子使劲往自己这边拽，于是氧端微微带负电，碳端微微带正电，成了一个“电子缺口”。

这个缺口就是羰基反应性的全部来源。带着富余电子的分子（亲核体）会被这个缺口吸引，扑上去形成新键——羰基碳是有机分子里最容易被“敲门”的位置之一。你身体此刻就在大规模利用这一点：葡萄糖分子里也有一个羰基，细胞的能量工厂从摆弄这个羰基开始拆解葡萄糖（第 46、52 章）。

醛和酮的差别只在位置上。醛的羰基在碳链头上，旁边还连着一个氢，这个氢很容易被氧“顺手牵羊”，于是醛特别容易被氧化成羧酸——乙醛变成乙酸，走的就是这条路。酮的羰基被夹在碳链中间，两边都是碳，没有那个顺手可牵的氢，所以酮安分得多。实验室里鉴别醛和酮的经典办法，正是利用这个差别：一种温和的氧化剂能啃动醛却啃不动酮，谁被氧化，谁就是醛。

生活里的羰基分子随处可见：洗甲水里那股挥发极快的气味来自丙酮（最简单的酮，溶解塑料和指甲油里的高分子是一把好手）；泡制生物标本的福尔马林是甲醛的水溶液；新装修房间里让人担忧的“异味”里，也常有甲醛这类小分子的身影。

\section{羧酸为什么酸，酯为什么香}

流水线第三站，羧酸。第 32 章说过，酸的本质是愿意放出氢离子（\chem{H^+}）的分子。可醇的羟基上也有氢，凭什么乙醇不显酸性，乙酸却酸得分明？

关键在“放手之后的日子好不好过”。一个分子放出了 \chem{H^+}，剩下的部分就得独自背负一个负电荷。负电荷越拥挤、越没处安放，这个分子就越不舍得放氢。乙醇放出氢之后，负电荷整个压在一个氧原子上，无处可逃——所以乙醇极少放氢，是中性的。而乙酸放出氢之后，负电荷由羧酸根里\textbf{两个氧原子共同分摊}（这就是第 41 章要细讲的“离域”：电子不困死在一个原子上）。负担减半，放氢就容易得多：

\[\chem{CH_3COOH} \rightleftharpoons \chem{CH_3COO^-} + \chem{H^+}\]

注意箭头是可逆的：乙酸是弱酸，每一百个分子里只有几个肯放氢，剩下的犹犹豫豫（第 31、32 章的平衡视角）。但这点酸性已经足够让味蕾尖叫，也足够让小苏打冒泡、让水垢溶解——你家厨房里发生的那些“醋的小魔术”，都是这几个氢离子干的。

现在，把这一章的两个角色请到一起：羧酸遇上醇，羧基和羟基各出一半，挤出一个水分子，剩下的部分手拉手连成酯：

\[\chem{CH_3COOH} + \chem{CH_3CH_2OH} \rightleftharpoons \chem{CH_3COOCH_2CH_3} + \chem{H_2O}\]

这个反应叫\textbf{酯化}，生成物乙酸乙酯带着清甜的果香。它还是个可逆反应（第 31 章）：常温下慢吞吞，两边僵持不下，工业上要靠酸催化加快脚步（催化剂只提速、不改终点，第 30 章），并设法把生成的水移走，才把平衡往酯的一边推。

酯类是水果香气的头号功臣。熟香蕉的甜香主要来自乙酸异戊酯，菠萝香里有丁酸乙酯，苹果和梨的香气里也各有一组酯的配方。想想这背后的演化逻辑：没熟的果实又酸又涩——种子还没长好，果实“拒绝”被吃；等种子成熟，果肉里的糖增多、酸下降、酯类香气分子大量合成，等于挂出一块“可以吃了”的广告牌，请动物们帮忙传播种子。你闻到的那缕香蕉味，是植物写了几亿年的广告文案（自然选择怎样写出这样的文案，第 76 章见）。

酯的香气分子通常不大、又不太会和水握手，所以容易挥发——这正是“香气能飘进鼻子”的原因。反过来，酯碰到水也可能慢慢水解，拆回羧酸和醇；有碱帮忙时，水解会快得多、也彻底得多。这最后一条性质，通向两样你每天都在用的东西。

\section{从肥皂到生物柴油}

第 47 章会正式介绍油脂：三个脂肪酸分子通过酯基挂在一个甘油骨架上。既然酯能水解，油脂当然也能。把油脂和烧碱（\chem{NaOH}）一起煮，酯键被剪断，甘油脱落，脂肪酸变成脂肪酸钠——这就是\textbf{肥皂}。这个反应有个专门的名字：皂化。

肥皂去污的本事，本章前面已经埋下了伏笔：脂肪酸钠分子一头是长长的碳链（不合群、亲油），一头是带电的羧酸根（会和水握手、亲水）。油污被碳链那头“抱住”，水分子拽着带电那头，把油污整团拖进水里——一个分子，同时是两边的“自己人”。这就是“相似相溶”规律的一次漂亮应用（第 25 章）。

同样的酯化学还开进了加油站。把植物油或餐饮废油与甲醇混合，在催化剂帮助下做“酯交换”——把甘油从酯基上换下来，换上小小的甲基——产物是脂肪酸甲酯，可以直接当柴油烧，这就是\textbf{生物柴油}。副产物甘油还能卖给日化厂。一辆车烧着炸过油条的地沟油跑上路，听起来像段子，背后却是实打实的酯交换反应：化学不管原料出身贵贱，只认结构。

\section{官能团说明书何时会失灵}

读到这里，你可能会觉得拿到了一本万能说明书：认出官能团，就能预言性格。是时候泼点冷水了——这本说明书有几条明确的“免责条款”。

第一，官能团的性格会被邻居改写。同样是羟基，挂在碳链上（乙醇）几乎不放氢，长在羧基里（乙酸）却肯放氢——因为隔壁的羰基在帮忙分摊负电荷。分子不是零件的简单拼装，零件之间会互相“说话”。

第二，大分子里接口太多，谁嗓门大要看场合。一个葡萄糖分子有五个羟基加一个羰基，它的行为是全体官能团的“委员会决议”，不是任何单一接口的写照（第 46 章）。

第三，气味这类感官性质尤其不听话。结构相似的酯可以有截然不同的气味，而结构相差很远的分子偶尔闻起来相似——嗅觉受体识别的是分子的三维形状拼图（第 42 章会讲形状多么重要），不是官能团名单。

第四，剂量与个体差异永远在场。本章反复出现的这条原则，对一切“某物质有没有害”的讨论都适用：离开剂量谈毒性，或者假设人人反应相同，都是把说明书当成了法律。

\section{回到那三瓶液体}

现在，兑现开场的承诺，把整条流水线画完整：

\begin{itemize}[itemsep=2pt]
\item \textbf{医用酒精}是乙醇，羟基让它与水互溶、易挥发、会降温；它进入身体后走上“乙醇→乙醛→乙酸”的拆解链。
\item \textbf{白醋}是乙酸的稀溶液，羧基让它肯放氢离子，于是酸得理直气壮；开瓶的酒变酸，就是这条链被醋酸菌走完的声音。
\item \textbf{香蕉的甜香}主要来自酯——羧酸和醇在果肉里握手、挤掉一分子水的产物；而熟过头的香蕉飘出酒味，是果肉细胞在缺氧时把糖酿成了乙醇。
\end{itemize}

看出来了吗？三样东西不在三条街上，而在同一条氧化链与酯化链的交叉路口：醇、醛（酮）、羧酸、酯，四种接口，两两之间都有路可走。更要命的一个发现是：这条流水线的上游，全都指向同一个源头——\textbf{糖}。酵母把糖酿成酒，细菌把酒变成醋，果实用糖造出酸和酯再引来动物。糖从哪里来？它为什么既能当燃料、又能当原料？那是下一篇的故事：第 46 章，我们去认识这种生命世界最忙碌的分子。

\begin{tiaozhan}
酸性的强弱可以用一个数字精确表达：$\mathrm{p}K_a$。它来自酸的电离平衡常数取负对数（第 31、32 章），数值越小，酸性越强。乙酸的 $\mathrm{p}K_a$ 约为 4.8，而乙醇的 $\mathrm{p}K_a$ 约为 16。注意对数的含义：相差约 11 个单位，意味着乙酸放氢的倾向约是乙醇的 $10^{11}$ 倍——一百亿倍！同是一个羟基氢，只因为隔壁多了一个羰基帮忙分摊负电荷，身价就差出一百亿倍。如果你想知道“离域”为什么能把负电荷摊平、共振结构式该怎么读，第 41 章会从苯环开始认真讲这件事。跳过本框不影响主线。
\end{tiaozhan}

\section*{练一练}

\begin{sikaoti}
\item 画出“乙醇→乙醛→乙酸”这条氧化链的示意图：用方框表示三种分子，框内标注各自的官能团（羟基、羰基、羧基），箭头旁注明每一步“多了一个氧还是少了一个氢”。图旁写一句：身体里的酶和醋缸里的醋酸菌，做的是同一种什么工作？
\item 1-丙醇（三个碳）和甘油（丙三醇，三个碳上各有一个羟基）相比，谁更易溶于水、谁更难挥发？用“羟基与碳链的拔河”说明理由，并预测它们沸点的高低顺序。
\item 把少量乙酸和小苏打（碳酸氢钠）混在一起会冒气泡。请写出乙酸的电离式，并说明：气泡里的气体是什么？这个实验能证明乙酸是强酸还是弱酸吗？为什么？
\item 画出肥皂分子的“物质流图”：一个长条分子，一端标“亲油”、一端标“亲水”，再画一小团油污和若干水分子，用箭头表示搅动后油污去了哪里。并回答：纯水和肥皂水都能冲掉灰尘，为什么只有肥皂水对付得了油污？
\item 乙酸乙酯可以水解回乙酸和乙醇。结合第 31 章的平衡观点回答：工厂想多生产酯时，为什么要设法把生成的水不断移走？在反应体系里加催化剂能不能提高酯的最终产量？
\item 两位同学喝了等量的酒，一位脸红，一位不红。用本章的代谢链解释两种现象的分子原因，并说明为什么“脸红说明更能喝”的说法恰好颠倒了因果。
\end{sikaoti}
